\documentclass[11pt,a4paper]{article}
\usepackage[utf8]{inputenc}
\usepackage[T1]{fontenc}
\usepackage[spanish,es-minimal,es-noshorthands]{babel}
\usepackage{lmodern}
\usepackage[sfdefault,lining]{FiraSans}
\usepackage[a4paper,top=2.3cm,bottom=2.2cm,left=1.8cm,right=1.8cm,headheight=1.3cm,headsep=0.4cm,footskip=1cm]{geometry}
\usepackage{amsmath,amssymb}
\usepackage{icomma}
\usepackage{graphicx}
\usepackage[table]{xcolor}
\usepackage{tikz}
\usepackage{fontawesome5}
\usepackage{ccicons}
\usepackage{fancyhdr}
\usepackage{lastpage}
\usepackage{tasks}
\usepackage{needspace}
\usepackage[most]{tcolorbox}
\usepackage{array}
\usepackage[hidelinks]{hyperref}
\AtBeginDocument{\hypersetup{pdftitle={Ficha de ejercicios de la unidad 4: Álgebra. Polinomios. 3.º de ESO},pdfsubject={Matemáticas 3.º ESO},pdfkeywords={lenguaje algebraico, monomios, polinomios, suma, resta, producto, identidades notables, factor común, división de polinomios, regla de Ruffini}}}

\definecolor{navy}{HTML}{1B2748}
\definecolor{accent}{HTML}{E8603A}
\definecolor{teal}{HTML}{1F9E90}
\definecolor{soft}{HTML}{EEF1F8}
\definecolor{muted}{HTML}{5D6685}

\setlength{\parindent}{0pt}
\setlength{\parskip}{3pt}
\settasks{label=\alph*), label-format=\bfseries\color{navy}, label-width=1.3em, item-indent=1.8em, column-sep=0.6em, after-item-skip=5pt, before-skip=3pt, after-skip=4pt}

% ---------- capçalera i peu ----------
\newcommand{\llicencia}{\href{https://creativecommons.org/licenses/by-nc-sa/4.0/deed.es}{\ccbyncsa}\ \ Licencia Creative Commons Atribución-NoComercial-CompartirIgual 4.0 Internacional (CC BY-NC-SA 4.0)}
\pagestyle{fancy}
\fancyhf{}
\fancyhead[L]{\raisebox{-0.35\height}{\includegraphics[height=0.95cm]{logo-gva.png}}}
\fancyhead[C]{\small\color{muted}Ficha de ejercicios · Unidad 4: Álgebra · 3.º de ESO}
\fancyhead[R]{\raisebox{-0.35\height}{\includegraphics[height=0.95cm]{logo-vermellar.png}}}
\fancyfoot[L]{\scriptsize\color{muted}\llicencia}
\fancyfoot[R]{\small\color{muted}\thepage}
\renewcommand{\headrulewidth}{0.4pt}
\renewcommand{\headrule}{\color{muted!40}\hrule height \headrulewidth}
\fancypagestyle{primera}{\fancyhf{}\renewcommand{\headrulewidth}{0pt}
  \fancyfoot[L]{\scriptsize\color{muted}\llicencia}
  \fancyfoot[R]{\small\color{muted}\thepage}}

% ---------- nivells i exercicis ----------
\newcommand{\estrelles}[1]{{\color{accent}\ifcase#1\or\faStar\faStar[regular]\faStar[regular]\or\faStar\faStar\faStar[regular]\or\faStar\faStar\faStar\fi}}
\newcounter{exer}
\newcommand{\ex}[1]{\par\needspace{4\baselineskip}\medskip\stepcounter{exer}\noindent{\color{navy}\bfseries\theexer.}\ \ {\small\estrelles{#1}}\ \ }

\newcommand{\seccio}[2]{\par\needspace{10\baselineskip}\bigskip
  \begin{tcolorbox}[colback=navy,colframe=navy,coltext=white,arc=2mm,boxrule=0pt,left=3mm,right=3mm,top=1.2mm,bottom=1.2mm]
  {\large\bfseries #1\quad #2}
  \end{tcolorbox}\vspace{-2pt}}

\newcolumntype{C}[1]{>{\centering\arraybackslash}p{#1}}

% quadre de recordatori (només a la fitxa d'exercicis)
\newcommand{\recorda}[1]{\par\smallskip\begin{tcolorbox}[colback=accent!7,colframe=accent!55,arc=2mm,boxrule=0.5pt,left=2.5mm,right=2.5mm,top=1mm,bottom=1mm]\small{\color{accent}\faLightbulb[regular]}\ \textbf{Recuerda:}\ #1\end{tcolorbox}}
\newcommand{\res}[1]{#1}

\def\fmneed{14}
% =====================================================================
% Capçalera FentMates de les fitxes: un color per a cada curs.
% Els colors són els de la pàgina inicial (LEVELS a assets/js/materials.js);
% el logo ve de _eines/logo/pdf/ (python3 _eines/logo/genera.py).
%
% Ús, al final del preàmbul de la unitat:
%   % =====================================================================
% Capçalera FentMates de les fitxes: un color per a cada curs.
% Els colors són els de la pàgina inicial (LEVELS a assets/js/materials.js);
% el logo ve de _eines/logo/pdf/ (python3 _eines/logo/genera.py).
%
% Ús, al final del preàmbul de la unitat:
%   % =====================================================================
% Capçalera FentMates de les fitxes: un color per a cada curs.
% Els colors són els de la pàgina inicial (LEVELS a assets/js/materials.js);
% el logo ve de _eines/logo/pdf/ (python3 _eines/logo/genera.py).
%
% Ús, al final del preàmbul de la unitat:
%   \input{../comu/fentmates.tex}
%   \fmcurs{1eso}
% i, al començament del document:
%   \capcalera{Matemàtiques}{1r ESO}{Fitxa d'exercicis · Unitat 2}{Divisibilitat}
% Opcional abans de \input: \def\fmneed{14} (espai mínim abans de cada secció).
% =====================================================================
\makeatletter
\def\fm@logodir{../../logo/pdf}
\providecommand{\fmneed}{10}
\def\fm@cursnom{}

% clau del curs, color i color del text damunt del color
\newcommand{\fm@curs}[3]{\@namedef{fm@color@#1}{#2}\@namedef{fm@text@#1}{#3}}
\fm@curs{1eso}{E6B84A}{navy}
\fm@curs{2eso}{3B82F6}{white}
\fm@curs{3eso}{FF7A50}{white}
\fm@curs{4eso-a}{A78BFA}{white}
\fm@curs{4eso-b}{6D28D9}{white}
\fm@curs{1batx-mi}{9B1C1C}{white}
\fm@curs{1batx-ccss}{EF6B6B}{white}
\fm@curs{2batx-mii}{1E6B47}{white}
\fm@curs{2batx-ccss}{2BB3A3}{white}

\newcommand{\fmcurs}[1]{%
  \@ifundefined{fm@color@#1}{\PackageError{fentmates}{Curs desconegut: #1}{}}{}%
  \edef\fm@tmp{\noexpand\definecolor{curs}{HTML}{\@nameuse{fm@color@#1}}%
               \noexpand\colorlet{curstext}{\@nameuse{fm@text@#1}}}\fm@tmp
  \colorlet{cursfosc}{curs!72!black}%
  \def\fm@logo{\fm@logodir/vertical-#1.pdf}%
  \def\fm@marca{\fm@logodir/marca-#1.pdf}}

% ---------- capçalera de la primera pàgina ----------
% #1 matèria (petit), #2 curs, #3 tipus de document i unitat, #4 títol
\newcommand{\capcalera}[4]{%
  \gdef\fm@cursnom{#2}%
  \par\noindent
  \begin{tcolorbox}[colback=curs,colframe=curs,boxrule=0pt,arc=2.2mm,width=3.5cm,height=2.15cm,
      valign=center,left=3mm,right=2mm,top=1mm,bottom=1mm,nobeforeafter,box align=center]
    \color{curstext}{\footnotesize #1}\par\vspace{1pt}{\fontsize{19}{21}\selectfont\bfseries #2\par}%
  \end{tcolorbox}\hspace{5mm}%
  \begin{minipage}[c]{\dimexpr\linewidth-3.5cm-5mm-2.9cm\relax}
    {\color{cursfosc}\small\bfseries\MakeUppercase{#3}}\par\vspace{2pt}
    {\color{navy}\fontsize{20}{23}\selectfont\bfseries #4\par}
  \end{minipage}\hfill
  \begin{minipage}[c]{2.6cm}\raggedleft\includegraphics[width=2.6cm]{\fm@logo}\end{minipage}\par
  \vspace{2pt}}

% ---------- barres de secció del color del curs ----------
\renewcommand{\seccio}[2]{\par\needspace{\fmneed\baselineskip}\bigskip
  \begin{tcolorbox}[colback=curs,colframe=curs,coltext=curstext,arc=2mm,boxrule=0pt,left=3mm,right=3mm,top=1.2mm,bottom=1.2mm]
  {\large\bfseries #1\quad #2}
  \end{tcolorbox}\vspace{-2pt}}

% ---------- estrelles de dificultat del color del curs (a Batxillerat no se'n posen) ----------
\@ifundefined{estrelles}{}{%
  \renewcommand{\estrelles}[1]{{\color{curs}\ifcase#1\or\faStar\faStar[regular]\faStar[regular]\or\faStar\faStar\faStar[regular]\or\faStar\faStar\faStar\fi}}}

% ---------- capçalera i peu de les pàgines ----------
% primera pàgina: logos institucionals al peu
\fancypagestyle{primera}{\fancyhf{}\renewcommand{\headrulewidth}{0pt}
  \fancyfoot[L]{\raisebox{-0.32\height}{\includegraphics[height=0.78cm]{logo-gva.png}}\hspace{5mm}%
    \raisebox{-0.32\height}{\includegraphics[height=0.78cm]{logo-vermellar.png}}}
  \fancyfoot[R]{\scriptsize\color{muted}\href{https://creativecommons.org/licenses/by-nc-sa/4.0/deed.ca}{\ccbyncsa}\ \ CC BY-NC-SA 4.0 · fentmates.org\qquad\small\thepage}}
% la resta: el símbol a l'esquerra i el curs a la dreta, amb el títol al centre com abans
\fancyhead[L]{\raisebox{-0.3\height}{\includegraphics[height=0.62cm]{\fm@marca}}}
\fancyhead[R]{\small\bfseries\color{cursfosc}\fm@cursnom}
\renewcommand{\headrule}{\color{curs!45}\hrule height \headrulewidth}
\makeatother

%   \fmcurs{1eso}
% i, al començament del document:
%   \capcalera{Matemàtiques}{1r ESO}{Fitxa d'exercicis · Unitat 2}{Divisibilitat}
% Opcional abans de \input: \def\fmneed{14} (espai mínim abans de cada secció).
% =====================================================================
\makeatletter
\def\fm@logodir{../../logo/pdf}
\providecommand{\fmneed}{10}
\def\fm@cursnom{}

% clau del curs, color i color del text damunt del color
\newcommand{\fm@curs}[3]{\@namedef{fm@color@#1}{#2}\@namedef{fm@text@#1}{#3}}
\fm@curs{1eso}{E6B84A}{navy}
\fm@curs{2eso}{3B82F6}{white}
\fm@curs{3eso}{FF7A50}{white}
\fm@curs{4eso-a}{A78BFA}{white}
\fm@curs{4eso-b}{6D28D9}{white}
\fm@curs{1batx-mi}{9B1C1C}{white}
\fm@curs{1batx-ccss}{EF6B6B}{white}
\fm@curs{2batx-mii}{1E6B47}{white}
\fm@curs{2batx-ccss}{2BB3A3}{white}

\newcommand{\fmcurs}[1]{%
  \@ifundefined{fm@color@#1}{\PackageError{fentmates}{Curs desconegut: #1}{}}{}%
  \edef\fm@tmp{\noexpand\definecolor{curs}{HTML}{\@nameuse{fm@color@#1}}%
               \noexpand\colorlet{curstext}{\@nameuse{fm@text@#1}}}\fm@tmp
  \colorlet{cursfosc}{curs!72!black}%
  \def\fm@logo{\fm@logodir/vertical-#1.pdf}%
  \def\fm@marca{\fm@logodir/marca-#1.pdf}}

% ---------- capçalera de la primera pàgina ----------
% #1 matèria (petit), #2 curs, #3 tipus de document i unitat, #4 títol
\newcommand{\capcalera}[4]{%
  \gdef\fm@cursnom{#2}%
  \par\noindent
  \begin{tcolorbox}[colback=curs,colframe=curs,boxrule=0pt,arc=2.2mm,width=3.5cm,height=2.15cm,
      valign=center,left=3mm,right=2mm,top=1mm,bottom=1mm,nobeforeafter,box align=center]
    \color{curstext}{\footnotesize #1}\par\vspace{1pt}{\fontsize{19}{21}\selectfont\bfseries #2\par}%
  \end{tcolorbox}\hspace{5mm}%
  \begin{minipage}[c]{\dimexpr\linewidth-3.5cm-5mm-2.9cm\relax}
    {\color{cursfosc}\small\bfseries\MakeUppercase{#3}}\par\vspace{2pt}
    {\color{navy}\fontsize{20}{23}\selectfont\bfseries #4\par}
  \end{minipage}\hfill
  \begin{minipage}[c]{2.6cm}\raggedleft\includegraphics[width=2.6cm]{\fm@logo}\end{minipage}\par
  \vspace{2pt}}

% ---------- barres de secció del color del curs ----------
\renewcommand{\seccio}[2]{\par\needspace{\fmneed\baselineskip}\bigskip
  \begin{tcolorbox}[colback=curs,colframe=curs,coltext=curstext,arc=2mm,boxrule=0pt,left=3mm,right=3mm,top=1.2mm,bottom=1.2mm]
  {\large\bfseries #1\quad #2}
  \end{tcolorbox}\vspace{-2pt}}

% ---------- estrelles de dificultat del color del curs (a Batxillerat no se'n posen) ----------
\@ifundefined{estrelles}{}{%
  \renewcommand{\estrelles}[1]{{\color{curs}\ifcase#1\or\faStar\faStar[regular]\faStar[regular]\or\faStar\faStar\faStar[regular]\or\faStar\faStar\faStar\fi}}}

% ---------- capçalera i peu de les pàgines ----------
% primera pàgina: logos institucionals al peu
\fancypagestyle{primera}{\fancyhf{}\renewcommand{\headrulewidth}{0pt}
  \fancyfoot[L]{\raisebox{-0.32\height}{\includegraphics[height=0.78cm]{logo-gva.png}}\hspace{5mm}%
    \raisebox{-0.32\height}{\includegraphics[height=0.78cm]{logo-vermellar.png}}}
  \fancyfoot[R]{\scriptsize\color{muted}\href{https://creativecommons.org/licenses/by-nc-sa/4.0/deed.ca}{\ccbyncsa}\ \ CC BY-NC-SA 4.0 · fentmates.org\qquad\small\thepage}}
% la resta: el símbol a l'esquerra i el curs a la dreta, amb el títol al centre com abans
\fancyhead[L]{\raisebox{-0.3\height}{\includegraphics[height=0.62cm]{\fm@marca}}}
\fancyhead[R]{\small\bfseries\color{cursfosc}\fm@cursnom}
\renewcommand{\headrule}{\color{curs!45}\hrule height \headrulewidth}
\makeatother

%   \fmcurs{1eso}
% i, al començament del document:
%   \capcalera{Matemàtiques}{1r ESO}{Fitxa d'exercicis · Unitat 2}{Divisibilitat}
% Opcional abans de \input: \def\fmneed{14} (espai mínim abans de cada secció).
% =====================================================================
\makeatletter
\def\fm@logodir{../../logo/pdf}
\providecommand{\fmneed}{10}
\def\fm@cursnom{}

% clau del curs, color i color del text damunt del color
\newcommand{\fm@curs}[3]{\@namedef{fm@color@#1}{#2}\@namedef{fm@text@#1}{#3}}
\fm@curs{1eso}{E6B84A}{navy}
\fm@curs{2eso}{3B82F6}{white}
\fm@curs{3eso}{FF7A50}{white}
\fm@curs{4eso-a}{A78BFA}{white}
\fm@curs{4eso-b}{6D28D9}{white}
\fm@curs{1batx-mi}{9B1C1C}{white}
\fm@curs{1batx-ccss}{EF6B6B}{white}
\fm@curs{2batx-mii}{1E6B47}{white}
\fm@curs{2batx-ccss}{2BB3A3}{white}

\newcommand{\fmcurs}[1]{%
  \@ifundefined{fm@color@#1}{\PackageError{fentmates}{Curs desconegut: #1}{}}{}%
  \edef\fm@tmp{\noexpand\definecolor{curs}{HTML}{\@nameuse{fm@color@#1}}%
               \noexpand\colorlet{curstext}{\@nameuse{fm@text@#1}}}\fm@tmp
  \colorlet{cursfosc}{curs!72!black}%
  \def\fm@logo{\fm@logodir/vertical-#1.pdf}%
  \def\fm@marca{\fm@logodir/marca-#1.pdf}}

% ---------- capçalera de la primera pàgina ----------
% #1 matèria (petit), #2 curs, #3 tipus de document i unitat, #4 títol
\newcommand{\capcalera}[4]{%
  \gdef\fm@cursnom{#2}%
  \par\noindent
  \begin{tcolorbox}[colback=curs,colframe=curs,boxrule=0pt,arc=2.2mm,width=3.5cm,height=2.15cm,
      valign=center,left=3mm,right=2mm,top=1mm,bottom=1mm,nobeforeafter,box align=center]
    \color{curstext}{\footnotesize #1}\par\vspace{1pt}{\fontsize{19}{21}\selectfont\bfseries #2\par}%
  \end{tcolorbox}\hspace{5mm}%
  \begin{minipage}[c]{\dimexpr\linewidth-3.5cm-5mm-2.9cm\relax}
    {\color{cursfosc}\small\bfseries\MakeUppercase{#3}}\par\vspace{2pt}
    {\color{navy}\fontsize{20}{23}\selectfont\bfseries #4\par}
  \end{minipage}\hfill
  \begin{minipage}[c]{2.6cm}\raggedleft\includegraphics[width=2.6cm]{\fm@logo}\end{minipage}\par
  \vspace{2pt}}

% ---------- barres de secció del color del curs ----------
\renewcommand{\seccio}[2]{\par\needspace{\fmneed\baselineskip}\bigskip
  \begin{tcolorbox}[colback=curs,colframe=curs,coltext=curstext,arc=2mm,boxrule=0pt,left=3mm,right=3mm,top=1.2mm,bottom=1.2mm]
  {\large\bfseries #1\quad #2}
  \end{tcolorbox}\vspace{-2pt}}

% ---------- estrelles de dificultat del color del curs (a Batxillerat no se'n posen) ----------
\@ifundefined{estrelles}{}{%
  \renewcommand{\estrelles}[1]{{\color{curs}\ifcase#1\or\faStar\faStar[regular]\faStar[regular]\or\faStar\faStar\faStar[regular]\or\faStar\faStar\faStar\fi}}}

% ---------- capçalera i peu de les pàgines ----------
% primera pàgina: logos institucionals al peu
\fancypagestyle{primera}{\fancyhf{}\renewcommand{\headrulewidth}{0pt}
  \fancyfoot[L]{\raisebox{-0.32\height}{\includegraphics[height=0.78cm]{logo-gva.png}}\hspace{5mm}%
    \raisebox{-0.32\height}{\includegraphics[height=0.78cm]{logo-vermellar.png}}}
  \fancyfoot[R]{\scriptsize\color{muted}\href{https://creativecommons.org/licenses/by-nc-sa/4.0/deed.ca}{\ccbyncsa}\ \ CC BY-NC-SA 4.0 · fentmates.org\qquad\small\thepage}}
% la resta: el símbol a l'esquerra i el curs a la dreta, amb el títol al centre com abans
\fancyhead[L]{\raisebox{-0.3\height}{\includegraphics[height=0.62cm]{\fm@marca}}}
\fancyhead[R]{\small\bfseries\color{cursfosc}\fm@cursnom}
\renewcommand{\headrule}{\color{curs!45}\hrule height \headrulewidth}
\makeatother

\fmcurs{3eso}

\begin{document}
\thispagestyle{primera}

\capcalera{Matemáticas}{3.º ESO}{Ficha de ejercicios · Unidad 4}{Álgebra. Polinomios}
\vspace{10pt}
Nombre y apellidos: \rule{6.6cm}{0.4pt}\hfill Grupo: \rule{1.6cm}{0.4pt}\hfill Fecha: \rule{2.2cm}{0.4pt}
\vspace{8pt}

\begin{tcolorbox}[colback=soft,colframe=soft,arc=2mm,left=3mm,right=3mm,top=1.5mm,bottom=1.5mm]
\small Ejercicios graduados por niveles de dificultad:\quad \estrelles{1}\ básico \qquad \estrelles{2}\ medio \qquad \estrelles{3}\ avanzado.\\[2pt] Resuélvelos en el cuaderno mostrando todos los pasos y da los polinomios ordenados de mayor a menor grado. Haz los cálculos sin calculadora.
\end{tcolorbox}

% =====================================================
\seccio{1}{Lenguaje algebraico}

\ex{1} Escribe en lenguaje algebraico:\begin{tasks}(2)
  \task El doble de un número más 5.
  \task El triple de un número menos su cuadrado.
  \task La mitad del siguiente de un número.
  \task El cuadrado de la suma de dos números.
  \task La suma de dos números consecutivos.
  \task La tercera parte de un número más su doble.
  \task El $15\,\%$ de un número.
  \task La diferencia de los cuadrados de dos números.
\end{tasks}\recorda{Fíjate en el orden: «el cuadrado de la suma» es $(x+y)^2$, pero «la suma de los cuadrados» es $x^2+y^2$.}

\ex{1} Escribe con palabras lo que expresa cada expresión algebraica:\begin{tasks}(4)
  \task $2x-1$
  \task $x^2+3$
  \task $\dfrac{x+y}{2}$
  \task $3(x-4)$
\end{tasks}

\ex{1} Calcula el valor numérico de cada expresión para los valores indicados:\begin{tasks}(2)
  \task $3x-5$, \ para $x=4$
  \task $x^2-2x$, \ para $x=-3$
  \task $-x^2+5$, \ para $x=-2$
  \task $2x^2y-xy$, \ para $x=1,\ y=-2$
\end{tasks}

\ex{2} Escribe una expresión algebraica para el perímetro y el área de cada figura:\begin{tasks}(2)
  \task Un rectángulo de base $3x$ y altura $x+2$.
  \task Un cuadrado de lado $x-1$.
\end{tasks}

\ex{2} Pedro tiene $x$ años; su hermana tiene el doble menos 3 y su padre, 25 más que Pedro.\begin{tasks}(1)
  \task Escribe la edad de cada uno y la suma de las tres edades.
  \task Escribe la suma de las tres edades dentro de 4 años.
  \task ¿Cuánto suman las tres edades si Pedro tiene 12 años?
\end{tasks}

\ex{2} Un aparcamiento cobra 1,80 € por entrar y 0,05 € por cada minuto. Escribe la expresión del precio de una estancia de $x$ minutos y calcula cuánto cuestan 2 horas.

\ex{3} Utiliza el lenguaje algebraico para demostrar que:\begin{tasks}(1)
  \task Si a un número le sumamos su doble, el resultado siempre es múltiplo de 3.
  \task La suma de dos números pares siempre es par.
  \task La suma de dos números impares consecutivos siempre es múltiplo de 4.
\end{tasks}


% =====================================================
\seccio{2}{Monomios}

\ex{1} Copia y completa la tabla:\par\smallskip{\centering\renewcommand{\arraystretch}{1.35}\begin{tabular}{|C{2.6cm}|C{2.6cm}|C{2.6cm}|C{1.6cm}|}\hline\rowcolor{soft}\textbf{Monomio} & \textbf{Coeficiente} & \textbf{Parte literal} & \textbf{Grado}\\\hline $-8x^3$ & & & \\\hline $a^2b$ & & & \\\hline $\tfrac{2}{3}xy^4$ & & & \\\hline $-y$ & & & \\\hline $12$ & & & \\\hline $5x^2y^3z$ & & & \\\hline \end{tabular}\par}\recorda{El \textbf{grado} de un monomio es la suma de los exponentes de la parte literal. Si no se escribe el coeficiente, vale 1.}

\ex{1} ¿Cuáles de estas expresiones son monomios? Justifica la respuesta.\begin{tasks}(4)
  \task $-3x^4$
  \task $\dfrac{5}{x}$
  \task $\dfrac{x^2}{4}$
  \task $2x+y$
  \task $\sqrt{x}$
  \task $7$
  \task $x^{-2}y$
  \task $-\tfrac12ab^2$
\end{tasks}

\ex{1} Agrupa los monomios semejantes:\quad $4x^2,\ \ -xy,\ \ 3x,\ \ -\tfrac12x^2,\ \ 5yx,\ \ x,\ \ 2x^2y,\ \ -x^2y,\ \ 9$.

\ex{2} Indica el grado de cada monomio respecto a cada variable y el grado total. Escribe, para cada uno, un monomio semejante.\begin{tasks}(3)
  \task $-6x^2y^5$
  \task $ab^3c$
  \task $\tfrac{3}{5}x^4y$
\end{tasks}

\ex{2} Calcula el valor numérico:\begin{tasks}(3)
  \task $-4x^3$ \ para $x=-2$
  \task $3xy^2$ \ para $x=2,\ y=-1$
  \task $\tfrac12x^2y$ \ para $x=4,\ y=3$
\end{tasks}

\ex{3} Resuelve:\begin{tasks}(1)
  \task Halla $k$ para que el monomio $-2x^{k}y^3$ sea de grado 8.
  \task Halla $m$ y $n$ para que $5x^{m+1}y^2$ y $-3x^4y^{n-1}$ sean semejantes.
  \task Escribe un monomio de grado 6 con tres variables, coeficiente $-\tfrac23$ y grado 3 respecto a $x$.
\end{tasks}


% =====================================================
\seccio{3}{Operaciones con monomios}

\ex{1} Calcula:\begin{tasks}(3)
  \task $3x+5x$
  \task $7x^{2}-2x^{2}$
  \task $-4ab+ab$
  \task $2x^{3}-x^{3}+5x^{3}$
  \task $xy-3xy-2xy$
  \task $6a^{2}b-6a^{2}b$
\end{tasks}\recorda{Solo se pueden sumar y restar monomios \textbf{semejantes}: sumamos los coeficientes y \textbf{no cambiamos} la parte literal.}

\ex{1} Reduce los términos semejantes:\begin{tasks}(2)
  \task $3x^{2}-5x+x^{2}+2x$
  \task $4xy-2x+3xy+7x$
  \task $2a^{2}-a+5-a^{2}+3a-8$
  \task $x^{3}-2x+3x^{3}+x^{2}-4x$
\end{tasks}

\ex{1} Multiplica:\begin{tasks}(3)
  \task $2x^{3}\cdot 5x^{2}$
  \task $(-3a^{2})\cdot (4a^{3})$
  \task $(6xy^{2})\cdot (-2x^{3}y)$
  \task $(\tfrac{1}{2}x^{4})\cdot (8x)$
  \task $(-x^{2}y)\cdot (-xy^{3})\cdot 2x$
  \task $(-5ab)\cdot (3a^{2}b^{2})$
\end{tasks}\recorda{Para multiplicar, multiplicamos los coeficientes y \textbf{sumamos los exponentes} de cada letra: $x^m\cdot x^n=x^{m+n}$.}

\ex{2} Divide. Indica si el resultado es un monomio.\begin{tasks}(3)
  \task $12x^5:(3x^2)$
  \task $(-20a^4b^3):(5a^2b)$
  \task $9x^3y^2:(-3x^3y)$
  \task $(-14x^6):(-7x^6)$
  \task $6x^2y:(4x)$
  \task $4x^2:(8x^3)$
\end{tasks}

\ex{2} Calcula las potencias:\begin{tasks}(4)
  \task $(2x^{3})^{2}$
  \task $(-3xy^{2})^{3}$
  \task $(-x^{2}y)^{4}$
  \task $(\tfrac{1}{2} x)^{2}$
\end{tasks}

\ex{2} Opera y simplifica (recuerda la jerarquía):\begin{tasks}(2)
  \task $3x\cdot2x^2-4x^3$
  \task $10x^5:(2x^2)+x\cdot x^2$
  \task $(2x)^3-5x\cdot x^2$
  \task $6a^2b:(3a)-ab\cdot4$
\end{tasks}

\ex{3} Halla el monomio que falta en cada caso:\begin{tasks}(1)
  \task $(\quad?\quad)\cdot3x^2=12x^5$
  \task $20x^4y^2:(\quad?\quad)=-5x^2y$
  \task $4x^3+(\quad?\quad)=-x^3$
  \task $(\quad?\quad)^2=9x^6y^2$ \ (dos soluciones)
\end{tasks}


% =====================================================
\seccio{4}{Polinomios}

\ex{1} Para cada polinomio: ordénalo e indica su grado, el número de términos, el coeficiente principal y el término independiente.\begin{tasks}(2)
  \task $P(x)=5x-3x^{4}+2-x^{2}$
  \task $P(x)=7-x^{3}$
  \task $P(x)=4x^{2}-3x+x^{5}$
  \task $P(x)=-2+x$
\end{tasks}\recorda{El \textbf{grado} de un polinomio es el grado del término de mayor grado. El \textbf{término independiente} es el que no tiene letra.}

\ex{1} ¿Qué polinomios del ejercicio anterior son completos? Escribe los incompletos en forma completa (con coeficientes 0).

\ex{1} Sea $P(x)=2x^3-x^2+3x-4$. Calcula:\begin{tasks}(4)
  \task $P(0)$
  \task $P(1)$
  \task $P(-2)$
  \task $P\!\left(\tfrac12\right)$
\end{tasks}

\ex{2} Indica el grado de cada polinomio:\begin{tasks}(3)
  \task $3x^2y^3-xy+5$
  \task $a^4-2a^2b^3+b$
  \task $x^2y^2-x^3+7xy$
\end{tasks}

\ex{2} Comprueba cuáles de los números $1,\ -1,\ 2,\ -2,\ 3$ y $-3$ son raíces del polinomio $P(x)=x^{3}-7x+6$.

\ex{3} Resuelve:\begin{tasks}(1)
  \task Halla $m$ para que $P(x)=2x^3-mx+4$ cumpla $P(-2)=0$.
  \task Halla $k$ para que el polinomio $Q(x)=x^2+kx-6$ tenga por raíz $x=3$.
  \task Escribe un polinomio de grado 4, incompleto, con coeficiente principal $-1$ y tal que $P(0)=5$.
\end{tasks}


% =====================================================
\seccio{5}{Suma y resta de polinomios}

\ex{1} Sean los polinomios $P(x)=3x^{3}-2x^{2}+x-5$, \ $Q(x)=-x^{3}+4x^{2}-6$ \ y \ $R(x)=2x^{2}-3x+1$. Calcula:\begin{tasks}(4)
  \task $P+Q$
  \task $P-Q$
  \task $Q+R$
  \task $R-P$
\end{tasks}\recorda{Para restar, sumamos el \textbf{opuesto}: cambiamos el signo de \textbf{todos} los términos del polinomio que restamos.}

\ex{1} Haz estas sumas y restas colocando los términos semejantes en columna:\begin{tasks}(1)
  \task $(4x^4-3x^2+2x)+(x^3+5x^2-2x+7)$
  \task $(x^3-6x+1)-(2x^3+x^2-6x-4)$
\end{tasks}

\ex{2} Con los polinomios del ejercicio 27, calcula:\begin{tasks}(3)
  \task $P-Q+R$
  \task $P-(Q-R)$
  \task $P-(Q+R)$
\end{tasks}

\ex{2} Con los polinomios del ejercicio 27, calcula $2P+3Q$ y $P-2R$.

\ex{2} Opera y reduce:\begin{tasks}(1)
  \task $(5x^2y-3xy+y^2)-(2x^2y+xy-4y^2)$
  \task $(a^2-2ab+b^2)+(a^2+2ab+b^2)-(2a^2-b^2)$
\end{tasks}

\ex{3} Halla el polinomio $X(x)$ en cada caso:\begin{tasks}(1)
  \task $P(x)+X(x)=Q(x)$, \ con $P$ y $Q$ los del ejercicio 27.
  \task $(2x^3+x-1)-X(x)=x^3-4x^2+3$
\end{tasks}


% =====================================================
\seccio{6}{Producto de polinomios}

\ex{1} Multiplica:\begin{tasks}(2)
  \task $3x(2x^{2}-5x+4)$
  \task $-2x^{2}(x^{3}-3x+1)$
  \task $4ab(a-2b)$
  \task $-x(x^{2}-x-1)$
\end{tasks}\recorda{Multiplicamos el monomio por \textbf{cada término} del polinomio (propiedad distributiva).}

\ex{1} Multiplica:\begin{tasks}(2)
  \task $(x+3)(x+5)$
  \task $(x-4)(x+2)$
  \task $(2x-1)(x-3)$
  \task $(3x+2)(2x-5)$
\end{tasks}

\ex{2} Multiplica:\begin{tasks}(1)
  \task $(x^{2}-3x+2)(x+4)$
  \task $(2x^{2}+x-1)(x-3)$
  \task $(x^{3}-2x+1)(x^{2}+x)$
  \task $(x^{2}+x+1)(x^{2}-x+1)$
\end{tasks}

\ex{2} Sean $P(x)=x^2-2x+3$ y $Q(x)=2x-1$. Calcula $P\cdot Q$ y $Q\cdot Q$. ¿Qué grado tiene cada producto?

\ex{2} Opera y reduce:\begin{tasks}(1)
  \task $2x(x-3)-(x+1)(x-2)$
  \task $(x+4)(x-1)-x(x+3)$
  \task $3(x^2-1)-(2x+1)(x-2)$
\end{tasks}

\ex{3} Calcula $(x-1)(x+2)(x+3)$ y $(x+1)(x-1)(x^2+1)$.

\ex{3} De una plancha rectangular de $(2x+3)$ cm por $(x+4)$ cm recortamos un cuadrado de lado $x$ cm. Escribe como polinomio el área que queda y calcúlala para $x=5$.


% =====================================================
\seccio{7}{Identidades notables}

\ex{1} Desarrolla:\begin{tasks}(3)
  \task $(x+3)^{2}$
  \task $(x-5)^{2}$
  \task $(x+4)(x-4)$
  \task $(a+1)^{2}$
  \task $(x-1)^{2}$
  \task $(y+7)(y-7)$
\end{tasks}\recorda{$(a+b)^2=a^2+2ab+b^2$, \quad $(a-b)^2=a^2-2ab+b^2$, \quad $(a+b)(a-b)=a^2-b^2$.}

\ex{2} Desarrolla:\begin{tasks}(3)
  \task $(2x+5)^{2}$
  \task $(3x-1)^{2}$
  \task $(2x+3y)(2x-3y)$
  \task $(x^{2}-2)^{2}$
  \task $(\tfrac{1}{2}x+2)^{2}$
  \task $(-x-3)^{2}$
\end{tasks}

\ex{2} Expresa como cuadrado de un binomio o como suma por diferencia:\begin{tasks}(3)
  \task $x^2+6x+9$
  \task $x^2-10x+25$
  \task $x^2-81$
  \task $4x^2+4x+1$
  \task $9x^2-16$
  \task $25x^2-30x+9$
\end{tasks}

\ex{2} Completa:\begin{tasks}(1)
  \task $(x+\ \square\ )^2=x^2+\ \square\ +49$
  \task $(\ \square\ -3)^2=4x^2-\ \square\ +9$
  \task $(x+\ \square\ )(x-\ \square\ )=x^2-64$
  \task $(3x+\ \square\ )^2=\ \square\ +12x+\ \square$
\end{tasks}

\ex{2} Opera y simplifica:\begin{tasks}(1)
  \task $(x+2)^2-(x-2)^2$
  \task $(x+3)(x-3)-(x-3)^2$
  \task $2(x+1)^2-(x+1)(x-1)$
\end{tasks}

\ex{3} Calcula mentalmente con las identidades notables:\begin{tasks}(4)
  \task $101^2$
  \task $99^2$
  \task $51\cdot49$
  \task $1002^2$
\end{tasks}

\ex{3} Demuestra que la diferencia entre los cuadrados de dos números consecutivos siempre es un número impar.


% =====================================================
\seccio{8}{Extracción de factor común}

\ex{1} Saca factor común:\begin{tasks}(3)
  \task $5x^{2}-15x$
  \task $8x^{3}+12x^{2}$
  \task $6x^{4}-9x^{3}+3x^{2}$
  \task $7x-14$
  \task $x^{5}-x^{3}+x^{2}$
  \task $10x^{3}-4x$
\end{tasks}\recorda{Factor común: el MCD de los coeficientes y las letras que están en \textbf{todos} los términos con el \textbf{menor} exponente. ¡Compruébalo multiplicando!}

\ex{1} Saca factor común:\begin{tasks}(2)
  \task $6x^{2}y-9xy^{2}$
  \task $4a^{3}b^{2}+2a^{2}b$
  \task $12x^{2}y^{3}-8xy^{2}+4xy$
  \task $15a^{2}-10ab+5a$
\end{tasks}

\ex{2} Saca factor común un binomio:\begin{tasks}(3)
  \task $x(x-2)+5(x-2)$
  \task $3(x+1)-x(x+1)$
  \task $2x(x+4)-(x+4)$
\end{tasks}

\ex{2} Factoriza al máximo: primero saca factor común y después aplica una identidad notable.\begin{tasks}(3)
  \task $x^{3}-4x$
  \task $2x^{2}-8x+8$
  \task $3x^{3}+6x^{2}+3x$
  \task $x^{4}-16x^{2}$
  \task $5x^{2}-45$
  \task $x^{3}-x$
\end{tasks}

\ex{3} Calcula de la manera más rápida sacando factor común:\begin{tasks}(3)
  \task $13\cdot47+13\cdot53$
  \task $98^2+2\cdot98$
  \task $37\cdot25+37\cdot75$
\end{tasks}

\ex{3} A un rectángulo de lados $x$ y $x+3$ le quitamos un rectángulo de lados $3$ y $x+3$. Escribe el área que queda, saca factor común y comprueba que es una diferencia de cuadrados.


% =====================================================
\seccio{9}{División de polinomios}

\ex{1} Haz estas divisiones e indica el cociente y el resto:\begin{tasks}(3)
  \task $(x^{2}+7x+10):(x+2)$
  \task $(x^{2}-3x+5):(x-1)$
  \task $(2x^{2}+x-3):(x-1)$
\end{tasks}\recorda{Seguimos dividiendo hasta que el grado del resto sea \textbf{menor} que el del divisor. Comprobación: $D=d\cdot C+R$.}

\ex{1} Divide. Cuidado con los términos que faltan en el dividendo:\begin{tasks}(2)
  \task $(x^{3}-2x^{2}+3x-1):(x-2)$
  \task $(3x^{3}+x-4):(x+1)$
\end{tasks}

\ex{2} Divide:\begin{tasks}(1)
  \task $(x^{4}-2x^{3}+x-3):(x^{2}+1)$
  \task $(2x^{4}+x^{3}-3x^{2}+x):(x^{2}-x+2)$
\end{tasks}

\ex{2} Divide $(6x^3-x^2+4x-5):(2x-1)$ y comprueba el resultado con la prueba de la división.

\ex{2} En una división de polinomios, el divisor es $x^2-3$, el cociente $2x+1$ y el resto $x-4$. ¿Cuál es el dividendo?

\ex{3} Al dividir $D(x)=x^3+x^2-3x+5$ entre un polinomio $d(x)$ obtenemos de cociente $x+2$ y de resto $-2x+3$. Halla $d(x)$.


% =====================================================
\seccio{10}{Regla de Ruffini}

\ex{1} Divide aplicando la regla de Ruffini:\begin{tasks}(2)
  \task $(x^{3}-3x^{2}+2x-4):(x-2)$
  \task $(x^{2}+6x+5):(x+1)$
  \task $(2x^{3}-5x^{2}+x+3):(x-3)$
\end{tasks}\recorda{Escribe los coeficientes del dividendo \textbf{ordenado y completo}. Si el divisor es $x-a$, a la izquierda ponemos $a$ (para $x+2$ ponemos $-2$).}

\ex{1} Divide por Ruffini. Cuidado con los términos que faltan:\begin{tasks}(2)
  \task $(x^{4}-3x^{2}+2x-1):(x+2)$
  \task $(3x^{3}-4x+5):(x-1)$
\end{tasks}

\ex{2} Divide por Ruffini:\begin{tasks}(2)
  \task $(x^{4}-16):(x-2)$
  \task $(x^{5}+1):(x+1)$
\end{tasks}

\ex{2} Calcula, sin hacer la división, el resto de:\begin{tasks}(3)
  \task $(x^3-5x+7):(x+3)$
  \task $(2x^4-x^2+1):(x-1)$
  \task $(x^{5}-2x^3+x):(x+2)$
\end{tasks}

\ex{2} Comprueba, con el teorema del resto, cuáles de los números $1,\ -1,\ 2,\ -2,\ 3,\ -3$ son raíces de $P(x)=x^{3}-3x^{2}-4x+12$.

\ex{3} Resuelve:\begin{tasks}(1)
  \task Halla $k$ para que la división $(2x^3+kx^2-x+3):(x+1)$ sea exacta.
  \task Halla $m$ para que el resto de $(x^3-mx+4):(x-2)$ sea 4.
\end{tasks}

\ex{3} Comprueba que $x=1$ es una raíz de $P(x)=x^3-2x^2-x+2$. Divide por Ruffini entre $x-1$ y, aplicando Ruffini otra vez al cociente, escribe $P(x)$ como producto de tres factores de grado 1.

\vfill
\begin{center}\color{muted}\small\faCheckCircle[regular]\ \ ¡Ánimo! Si te cuesta algún ejercicio, repasa los apuntes interactivos de la unidad.\end{center}

\end{document}
